\section{热力学系统}

热学中把所研究的对象称为热力学系统，系统是由大量微观粒子所组成的宏观物体，系统之外与系统间发生相互作用的部分称为\uwave{外界}或\uwave{环境}。在一般的情况中，系统与外界之间总会存在某些相互作用，进而影响系统状态，系统可以依据其与外界相互作用的关系分为三种类型。

\subsection{热力学系统的分类}
\begin{BoxDefinition}[开放系统]
    \uwave{开放系统}，是指与外界既有能量交换，又有物质交换的系统。    
\end{BoxDefinition}

\begin{BoxDefinition}[封闭系统]
    \uwave{封闭系统}，是指与外界仅有能量交换，而无物质交换的系统。    
\end{BoxDefinition}

\begin{BoxDefinition}[孤立系统]
    \uwave{孤立系统}，是指与外界既无能量交换，又无物质交换的系统。    
\end{BoxDefinition}

例如，一杯敞开盖子的热水就是开放系统，因为水的温度再不断降低的同时还有水蒸气进入空气中。如果将盖子合上，此时就变为了封闭系统。如果将杯子进一步换成理想保温杯，此时就将变为孤立系统。孤立系统是比较理想的概念，在本章中，我们主要讨论的是封闭系统。\vspace{-0.2cm}

\subsection{热力学系统的参量}
在热力学系统中为了确定热力学系统的状态而引入的物理量，称为系统的\uwave{物态参量}。热学研究与力学研究不同，我们并不关心系统整体的宏观机械运动状态，例如位移、速度、加速度等现在并不是我们现在关注的物理量，我们将研究的注意力转向系统内部，组成系统的大量微观粒子处于永不停息的无规则运动中，每个粒子都具有质量、速度、动量、能量等，这些物理量统称为\uwave{微观量}，微观量在热学实验中一般不能直接观察和测量，而要通过统计物理的方法进行研究。而与之对应的，将描述热力学系统所表现的各种宏观性质的物态参量为\uwave{宏观量}。

在具体问题中如何选用物态参量，需要视具体情况而定，对于本章讨论的气体动理论，我们主要研究的是化学成分单一的气体系统，因此通常用\uwave{体积}和\uwave{压强}两个力学参量描述系统状态。

\empx{体积表征了气体分子所能达到的空间}，通常用$V$表示，需要注意的是，气体的体积与容器的容积不同，后者是气体活动空间和气体分子本身体积的和，只有在不计分子大小的情况下两者才能等效。体积的国际单位制是立方米（\si{m^3}），常用单位还有升（\si{L}），换算关系是
\begin{Equation}*
    1\si{m^3}=1000\si{L}=1000\si{dm^3}
\end{Equation}

\empx{压强表征了气体分子对容器壁的冲击}，通常用$P$表示，压强的国际单位制是帕斯卡（\si{Pa}），我们知道，压强衡量的是单位面积上的压力，因此 \si{Pa=N\cdot m^{-2}}。压强由于历史上的原因，现在仍然有许多常用的非国际单位，比较常用的单位还有标准大气压（\si{atm}）和毫米汞柱（\si{mmHg}）
\begin{Equation}*
    1\si{atm}=1.013\times 10^5\si{Pa}=760\si{mmHg}
\end{Equation}
这里的体积和压强均为力学参量，而在热学中我们还需要引入一个额外的热学参量，用来表征系统的热学性质，或者就是通俗所说的冷热程度，即温度。但到底什么是温度呢？为了解决这个问题，我们将在下面两小节中通过热平衡的概念和热平衡定律，引入温度的科学定义。

\subsection{热平衡定律}
热力学系统在不受外界影响的条件下，系统各部分的宏观性质不随时间变化的状态，我们称之为\uwave{平衡态}，反之称为\uwave{非平衡态}，深入研究表明，经过适当的时间，偏离平衡态不远的系统总是可以达到平衡态，由非平衡态达到平衡态所经历的时间，称为热力学系统的\uwave{弛豫时间}。

热力学系统处于平衡态时，不随时间变化的是所有能观测到的系统的宏观性质，而在微观上，组成系统的大量微观粒子仍然处于无规则的热运动中，只不过，这种热运动在宏观上表现出的平均效果不再随时间而变。因此，我们将这种平衡中蕴含着运动的状态，称为\uwave{热动平衡}。

假设有两个系统通过导热壁相互接触后达到一个相同的平衡态，我们就称这两个系统达到了\uwave{热平衡}。现在如\xref{fig:热平衡定律}所示，设想三个热力学系统A,B,C，置于一个孤立的环境中，系统A,B用绝热壁隔开，但使它们与处于确定状态的系统C通过导热壁接触，我们知道，经过足够长的时间后，系统A,B分别都将和C达到热平衡，而此时，如果将导热壁和绝热壁交换，我们将意外的发现系统A,B的状态并没有发生任何变化，这就表明系统A,B已经达到热平衡了。

\begin{Figure}[热平衡定律]
    \begin{FigureSub}
        \includegraphics{build/Chapter01A_01.fig.pdf}
    \end{FigureSub}\hspace{1cm}
    \begin{FigureSub}
        \includegraphics{build/Chapter01A_02.fig.pdf}
    \end{FigureSub}
\end{Figure}

\uwave{热平衡定律}就是对上述实验结果的概况
\begin{BoxLaw}[热平衡定律]
    在与外界隔绝的条件下，如果两个系统分别与第三个处于确定状态下的系统处于热平衡

    那么，这两个系统彼此之间也必定处于热平衡。
\end{BoxLaw}
热平衡定律是由福勒在1930年前后提出的，而在此之前，热力学第一定律和热力学第二定律都已确立，为了想说明其在逻辑上先于第一第二定律，也将热平衡定律称为\uwave{热力学第零定律}。

热平衡定律揭示了一切互为热平衡的系统之间必定存在一个共同的宏观热学性质。更深入的说，每个热力学系统都有一个特别的热学状态函数，当几个系统间达成热平衡时，它们各自的这一状态函数就取相同值，而热平衡定律保证了这种规定的可行性：系统A,B分别与C达成热平衡时，系统A,B间也是热平衡的，这样就可以确保热学状态函数具有等号的传递性。

而上述的这个热学状态函数，我们就定义为\uwave{温度}，这就是温度的科学定义。


为了定量的确定温度的数值，我们还必须给出温度的数值表示，即温标。我们在理论中比较常用的是\uwave{热力学温标}，记作$T$，热力学温标是基于热学现象引入的温标，或者说，热力学温标可以使得热学定律的数学形式变得简单，其也是温度的国际单位制，单位是开尔文（\si{K}）。另外在生活中比较常用的还有\uwave{摄氏温标}，记作$t$，单位是摄氏度（\si{\degreeCelsius}）。两种温标的换算关系是
\begin{Equation}*
    T(\si{K})=t(\si{\degreeCelsius})+273.15
\end{Equation}
即摄氏温标和热力学温标对于温度变化时的温标增量不同，但温度计量的起始基点不同。